\appendix

\section{Implementation details}
\label{app:impl}

\paragraph{Model.}
All experiments use Janus-Pro-7B (\texttt{deepseek-community/Janus-Pro-7B}),
loaded in \texttt{bfloat16} via the Hugging Face
\texttt{JanusForConditionalGeneration} class. Janus-Pro is a unified
autoregressive model with a shared transformer backbone and \emph{decoupled}
modality encoders: a SigLIP vision encoder for understanding (image$\to$text) and
a separate vector-quantized (VQ) tokenizer with an image head for generation
(text$\to$image), where an image is a $24\times24=576$ grid of discrete tokens
from a codebook of size $16384$. We use the two directions through the
processor's text mode (text language modelling) and image mode
(\texttt{max\_new\_tokens}$=576$ then VQ decoding). LoRA is applied \emph{only} to
the shared backbone; the SigLIP and VQ encoders/heads are frozen, which is what
makes the modality-localized weight analysis of Appendix~\ref{app:dw}
meaningful.

\paragraph{LoRA configuration (fixed across all runs).}
Every adapter targets the seven projections of each backbone layer---%
\texttt{q,k,v,o\_proj} (attention) and \texttt{gate,up,down\_proj} (MLP), with
$\alpha=2r$ (so the LoRA scaling $\alpha/r=2$ is constant across the rank sweep),
\texttt{bias=none}, causal-LM task type, and no rsLoRA; the induced update is
$\Delta W=(\alpha/r)BA$. Only the swept variable (rank $r$ or KL weight $\lambda$)
changes within a sweep.

\paragraph{Optimization (common).}
AdamW with $\beta=(0.9,0.95)$ (matching the Janus pre-training $\beta_2$) and zero
weight decay, a cosine schedule with $5\%$ warmup, gradient accumulation $8$,
gradient clipping $1.0$, \texttt{bfloat16}, single H100. With the regularizer the
loss is
$\mathcal{L}=\mathrm{CE}+\lambda\,\mathrm{KL}(p_{\text{adapter}}\|p_{\text{base}})$,
where $p_{\text{base}}$ is a \texttt{disable\_adapter()} no-grad forward on the
same inputs ($\approx$$2\times$ forward cost).

\paragraph{Uncertainty quantification.}
We report each number as $\text{mean}\pm$ one standard error, computed without
retraining and with an estimator-appropriate method per metric type. The accuracy
benchmarks (MMLU, POPE, MME) are proportions of correct items over a fixed
evaluation set, so their standard error follows the binomial/Wilson form
\citep{wilson1927,brown2001interval}. Per-instance metrics averaged over the
generated set (CLIP directional similarity, DINO similarity, CLIP-IQA) use
$s/\sqrt{n}$ over the $n$ generations, equivalently a nonparametric bootstrap
\citep{efron1993bootstrap}. The set-level distributional metrics (CMMD, KID) have
no closed-form variance, so their standard error is the standard deviation of a
nonparametric bootstrap over the generated and reference sets ($B{=}2000$
resamples) \citep{efron1993bootstrap,binkowski2018demystifying}. Perplexity is a
monotone transform of the mean per-chunk negative log-likelihood, whose standard
error propagates through the exponential. All standard errors quantify
evaluation-set sampling noise at fixed model weights; they do not capture
training-seed variance. The small benchmark movements we observe (e.g.\ MMLU
$0.492\!\to\!0.479$, within $\sim\!1$ standard error) thus lie within sampling
noise, whereas the acquisition shifts (e.g.\ CMMD, DirSim) do not. This is the
basis for the cautious reading of the boundary settings in
Section~\ref{sec:boundary}.

\paragraph{Text (poetry) runs.}
Data: \texttt{biglam/gutenberg-poetry-corpus}, grouped by source document with
front-matter removed and windowed to $512$ tokens ($6000$ train / $200$ val
windows). Fixed: $r{=}16$, $\alpha{=}32$, dropout $0.05$, LR $10^{-4}$, $3$
epochs ($150$-min wall-clock budget). Rank sweep: $r\in\{4,8,16,32,64\}$ at
$\lambda{=}0.5$; $\lambda$ sweep: $\lambda\in\{0,0.25,0.5,1,1.5\}$ at $r{=}16$.

\paragraph{Image (Naruto) runs.}
Data: \texttt{lambdalabs/naruto-blip-captions} ($\approx$$1141$ image--caption
pairs); acquisition on $80$ held-out prompts. Final configuration
(Table~\ref{tab:panel}): $r{=}48$, $\alpha{=}96$, dropout $0.10$, LR $10^{-4}$,
$6$ epochs, no caption prefix, $\lambda{=}0.5$. Rank sweep:
$r\in\{2,\dots,64\}$ at $\lambda{=}0.5$; $\lambda$ sweep at $r{=}32$.

\paragraph{Medicine, Welsh and sketch runs.}
Medicine: \texttt{ccdv/pubmed-summarization}, $r{=}32$, $\alpha{=}64$,
$\lambda{=}0.25$. Welsh: \texttt{wikimedia/wikipedia} (Welsh subset), $r{=}32$,
$\alpha{=}64$, $\lambda{=}0.1$, $3$ epochs. Sketch: \texttt{zoheb/sketch-scene},
identical configuration to the Naruto Final run ($r{=}48$, $\alpha{=}96$,
$\lambda{=}0.5$, LR $10^{-4}$, $6$ epochs). All three share the common
optimization settings above.
% TODO FC: add LR / epochs / corpus size for medicine, and MAX_DOCS /
% TRAIN_WINDOWS for Welsh, from the SLURM run scripts.

\paragraph{Prompts and image generation.}
Janus generates images autoregressively from a natural-language prompt
(text$\to$image); every generated image in this paper is therefore conditioned on
an explicit text prompt. During training the conditioning text is the dataset's
BLIP caption; the \textit{Capacity} and \textit{+KL} configurations prepend a
fixed style trigger (``anime sprite in Naruto style, ''), which \textit{Scaled}
and \textit{Final} remove, so the Final adapter applies the style
\emph{unconditionally}, with no trigger token at inference. Acquisition metrics
(Table~\ref{tab:panel}) are computed on images generated from $80$ held-out Naruto
captions. The qualitative grids instead use a fixed, curated set of generic
natural-language prompts mixing in-domain and clearly out-of-domain subjects
(e.g.\ ``a young man with blonde spiky hair'', ``a samurai warrior holding a
katana'', ``an old wizard with a long white beard''), chosen to test whether the
style transfers to arbitrary content; the exact prompt is printed beneath each row
of every grid (Figures~\ref{fig:janus_qualitative_comparison},
\ref{fig:rank_visual}, \ref{fig:kl_visual}). The cross-modal check under the text
adapter (Figure~\ref{fig:t2i_text_kl}) uses four neutral prompts (e.g.\ ``a cat
sitting on a chair''). All conditions share the same seed schedule, so within a
row the adapter is the only source of variation.


\section{Selecting the language target}
\label{app:langprobe}
Before training anything, we scored candidate language targets by
$\times$control: the base model's perplexity on the candidate corpus divided by
its perplexity on an English WikiText control (control perplexity $14.22$ under
this probe's chunking). Candidates were real, available corpora: low-resource
natural languages via Wikipedia, and a synthetic substitution-cipher language
built from English as a maximally distant anchor. Table~\ref{tab:langprobe}
reports the result.

We did not take the most distant candidate. The cipher is by far the furthest
($10.3\times$), but it is a deterministic character-level remapping of English:
an adapter can fit it by learning the substitution table, which would make the
resulting ``acquisition'' a shortcut rather than a new linguistic capability.
Klingon is marginally further than Welsh ($2.51\times$ vs.\ $2.34\times$) but its
available corpus is a few thousand short sentences, too small to separate
acquisition from memorization at our training budget. Welsh combines a genuine
distance from the base with an abundant, naturally occurring corpus, which is why
we selected it.

\begin{table}[h]\centering\small
\setlength{\tabcolsep}{8pt}
\begin{tabular}{lccl}
\toprule
Candidate & Base PPL & $\times$control & Note \\
\midrule
Cipher (substitution) & $146.8$ & $10.33$ & deterministic remapping of English \\
Klingon               & $35.7$  & $2.51$  & constructed; corpus too small \\
\textbf{Welsh}        & $33.3$  & $\mathbf{2.34}$ & \textbf{selected}: distant and abundant \\
Yoruba                & $24.0$  & $1.69$  & closer to the base \\
\bottomrule
\end{tabular}
\caption{Language-target probe, base model only, no training. $\times$control is
the base perplexity on the candidate divided by the base perplexity on an English
control ($14.22$ here), the same construction used for the poetry and medicine
targets. Higher = the base is more surprised by the corpus. We selected Welsh
rather than the two more distant candidates for the reasons given in the text.}
\label{tab:langprobe}
\end{table}


\section{Image-generation configuration: incremental ablation}
\label{app:ablation}
We reach the Final configuration by adding components one or a few at a time from
a naive starting point (attention-only LoRA, rank 8, with a fixed caption-prefix
trigger), which degraded the base and is omitted from Table~\ref{tab:panel}.
\textit{Capacity} restores trainable capacity (rank 16, LoRA on attention and
MLP); \textit{+KL} adds the KL-to-base regularizer ($\lambda{=}0.5$);
\textit{Scaled} removes the training-time trigger prefix and raises rank, learning
rate and dropout (rank 32, $8\times10^{-5}$, dropout 0.1, six epochs);
\textit{Final} raises the rank to 48 and the learning rate to $10^{-4}$. The
\textit{Joint} conditions additionally alternate a $50/50$ domain-controlled text
stream (matched: Naruto franchise; far: physics and mathematics) into the same
adapter. Table~\ref{tab:panel} reports the acquisition panel; the corresponding
generations are in Figure~\ref{fig:janus_qualitative_comparison}.

\begin{table}[t]
\centering\small
\setlength{\tabcolsep}{5pt}
\renewcommand{\arraystretch}{1.05}
\begin{tabular}{lcccccc}
\toprule
 & \textbf{Capacity} & \textbf{+KL} & \textbf{Scaled} & \textbf{Final} & \textbf{Joint} & \textbf{Joint} \\
 & & & & & \textbf{(matched)} & \textbf{(far)} \\
\midrule
\multicolumn{7}{l}{\textit{Configuration (bold = changed vs.\ previous column)}}\\
Rank $r$              & 16 & 16 & \textbf{32} & \textbf{48} & 48 & 48 \\
LR ($\times10^{-5}$)  & 5 & 5 & \textbf{8} & \textbf{10} & 10 & 10 \\
Epochs               & 8 & 8 & \textbf{6} & 6 & 6 & 6 \\
LoRA dropout         & .05 & .05 & \textbf{.10} & .10 & .10 & .10 \\
Train prefix         & yes & yes & \textbf{no} & no & no & no \\
KL $\lambda$         & 0 & \textbf{0.5} & 0.5 & 0.5 & 0.5 & 0.5 \\
Joint text           & --- & --- & --- & --- & \textbf{matched} & \textbf{far} \\
\midrule
\multicolumn{7}{l}{\textit{Acquisition (vs.\ real style, $N{=}80$; bold = best)}}\\
KID $\downarrow$       & $.039$ & $.017$ & $.016$ & $.016$ & $.016$ & $\mathbf{.014}$ \\
CMMD $\downarrow$      & $.189$ & $\mathbf{.105}$ & $.120$ & $.110$ & $.131$ & $.134$ \\
DINO sim $\uparrow$    & $.629$ & $\mathbf{.663}$ & $.624$ & $.640$ & $.633$ & $.637$ \\
CLIP-IQA $\uparrow$    & $.659$ & $\mathbf{.750}$ & $.694$ & $.701$ & $.725$ & $.704$ \\
CLIP DirSim $\uparrow$ & $-.059$ & $-.033$ & $.157$ & $\mathbf{.163}$ & $.146$ & $.151$ \\
\bottomrule
\end{tabular}
\caption{Image-generation configuration as an incremental ablation: each column
adds one or more components to the previous one to reach \textit{Final}.
\textit{Capacity}: rank 16 with LoRA on attention and MLP. \textit{+KL}: adds the
KL-to-base regularizer. \textit{Scaled}: removes the training-time caption prefix
and raises rank, learning rate and dropout. \textit{Final}: rank 48, LR
$10^{-4}$. \textit{Joint}: alternates a $50/50$ domain-controlled text stream
(matched: Naruto franchise; far: physics/mathematics) into the same adapter. The
naive start (attention-only, rank 8, prefix-gated) degraded the base and is
omitted; CLIP DirSim turning positive at \textit{Scaled} marks the onset of
genuine style transfer.}
\label{tab:panel}
\end{table}

\paragraph{The anchor strength has an acquisition optimum.}
Table~\ref{tab:klimg} sweeps $\lambda$ at rank 32 in the image direction. Unlike
the text direction, acquisition here is a single-peaked function of $\lambda$: it
peaks near $\lambda{=}0.25$ to $0.5$ and collapses for $\lambda\!\geq\!1$, where
the adapter is over-constrained toward the base. $\lambda{=}0$ is already
workable, so in this direction the anchor refines rather than enables the recipe.
Retention stays at the base level across the whole $\lambda$ range, so the anchor
changes acquisition without any measurable cross-capability cost. The CLIP style
metrics place the optimum at $\lambda{=}0.25$, yet we adopt $\lambda{=}0.5$ for
the Final recipe because those metrics reward proximity to the style
\emph{distribution} but not whether the prompt's content survives the
stylisation; Appendix~\ref{app:metricdebug} settles this with two content-aware
embedding measures and an independent VLM judge.

\begin{table}[t]
\centering\small
\setlength{\tabcolsep}{6pt}
\begin{tabular}{lcccccc}
\toprule
Metric & Base & $\lambda{=}0$ & $\lambda{=}0.25$ & $\lambda{=}0.5$ & $\lambda{=}1.0$ & $\lambda{=}1.5$ \\
\midrule
\multicolumn{7}{l}{\textit{Acquisition (vs.\ real style images)}}\\
CMMD $\downarrow$      & --- & $.155$ & $\mathbf{.100}$ & $.108$ & $.205$ & $.221$ \\
DINO sim $\uparrow$    & --- & $.626$ & $\mathbf{.647}$ & $.643$ & $.439$ & $.388$ \\
CLIP DirSim $\uparrow$ & --- & $.157$ & $\mathbf{.164}$ & $.152$ & $.049$ & $.024$ \\
CLIP-IQA $\uparrow$    & --- & $.666$ & $.691$ & $.691$ & $.738$ & $\mathbf{.739}$ \\
\midrule
\multicolumn{7}{l}{\textit{Retention}}\\
WikiText PPL $\downarrow$ & $13.71$ & $13.67$ & $13.68$ & $13.78$ & $13.76$ & $13.76$ \\
MMLU $\uparrow$           & $.492$ & $.495$ & $.493$ & $.497$ & $.493$ & $.492$ \\
POPE acc $\uparrow$       & $.877$ & $.879$ & $.877$ & $.879$ & $.877$ & $.873$ \\
MME $\uparrow$            & $1785$ & $1783$ & $1776$ & $1796$ & $1802$ & $1810$ \\
\bottomrule
\end{tabular}
\caption{Image-generation $\lambda$ sweep at rank 32. \textit{Acquisition} (bold
= best) peaks near $\lambda{=}0.25$ to $0.5$ and degrades for
$\lambda\!\geq\!1$. \textit{Retention} stays at the base level across the whole
$\lambda$ range (same probe configuration as Table~\ref{tab:rank}), so the anchor
changes acquisition without any measurable cross-capability cost.}
\label{tab:klimg}
\end{table}

\begin{figure*}[p]
\centering
\setlength{\tabcolsep}{1.5pt}
\renewcommand{\arraystretch}{1.0}
\small
\begin{tabular}{ccccccc}
\textbf{No adapter} & \textbf{Capacity} & \textbf{+KL} & \textbf{Scaled} & \textbf{Final} & \textbf{Joint (matched)} & \textbf{Joint (far)} \\
\textit{base 7B} & \textit{r16, +MLP} & \textit{+ KL reg} & \textit{r32, untriggered} & \textit{r48, LR $10^{-4}$} & \textit{+ Naruto wiki} & \textit{+ physics wiki} \\
\midrule
\includegraphics[width=0.135\textwidth]{figs/janus_comparison/no_adapter/eval_00_natural.png} &
\includegraphics[width=0.135\textwidth]{figs/janus_comparison/adapter_v2/eval_00_natural.png} &
\includegraphics[width=0.135\textwidth]{figs/janus_comparison/adapter_v3/eval_00_natural.png} &
\includegraphics[width=0.135\textwidth]{figs/janus_comparison/adapter_v4/eval_00_natural.png} &
\includegraphics[width=0.135\textwidth]{figs/janus_comparison/adapter_solo_gen/eval_00_natural.png} &
\includegraphics[width=0.135\textwidth]{figs/janus_comparison/adapter_text_matched/eval_00_natural.png} &
\includegraphics[width=0.135\textwidth]{figs/janus_comparison/adapter_text_far/eval_00_natural.png} \\
\multicolumn{7}{l}{\footnotesize\textit{Prompt (eval\_00):} a red panda with a black face and brown eyes} \\
\addlinespace[2pt]
\includegraphics[width=0.135\textwidth]{figs/janus_comparison/no_adapter/eval_01_natural.png} &
\includegraphics[width=0.135\textwidth]{figs/janus_comparison/adapter_v2/eval_01_natural.png} &
\includegraphics[width=0.135\textwidth]{figs/janus_comparison/adapter_v3/eval_01_natural.png} &
\includegraphics[width=0.135\textwidth]{figs/janus_comparison/adapter_v4/eval_01_natural.png} &
\includegraphics[width=0.135\textwidth]{figs/janus_comparison/adapter_solo_gen/eval_01_natural.png} &
\includegraphics[width=0.135\textwidth]{figs/janus_comparison/adapter_text_matched/eval_01_natural.png} &
\includegraphics[width=0.135\textwidth]{figs/janus_comparison/adapter_text_far/eval_01_natural.png} \\
\multicolumn{7}{l}{\footnotesize\textit{Prompt (eval\_01):} a man with glasses and a hoodie} \\
\addlinespace[2pt]
\includegraphics[width=0.135\textwidth]{figs/janus_comparison/no_adapter/eval_02_natural.png} &
\includegraphics[width=0.135\textwidth]{figs/janus_comparison/adapter_v2/eval_02_natural.png} &
\includegraphics[width=0.135\textwidth]{figs/janus_comparison/adapter_v3/eval_02_natural.png} &
\includegraphics[width=0.135\textwidth]{figs/janus_comparison/adapter_v4/eval_02_natural.png} &
\includegraphics[width=0.135\textwidth]{figs/janus_comparison/adapter_solo_gen/eval_02_natural.png} &
\includegraphics[width=0.135\textwidth]{figs/janus_comparison/adapter_text_matched/eval_02_natural.png} &
\includegraphics[width=0.135\textwidth]{figs/janus_comparison/adapter_text_far/eval_02_natural.png} \\
\multicolumn{7}{l}{\footnotesize\textit{Prompt (eval\_02):} a person with black hair and a green shirt} \\
\addlinespace[2pt]
\includegraphics[width=0.135\textwidth]{figs/janus_comparison/no_adapter/human_05_natural.png} &
\includegraphics[width=0.135\textwidth]{figs/janus_comparison/adapter_v2/human_05_natural.png} &
\includegraphics[width=0.135\textwidth]{figs/janus_comparison/adapter_v3/human_05_natural.png} &
\includegraphics[width=0.135\textwidth]{figs/janus_comparison/adapter_v4/human_05_natural.png} &
\includegraphics[width=0.135\textwidth]{figs/janus_comparison/adapter_solo_gen/human_05_natural.png} &
\includegraphics[width=0.135\textwidth]{figs/janus_comparison/adapter_text_matched/human_05_natural.png} &
\includegraphics[width=0.135\textwidth]{figs/janus_comparison/adapter_text_far/human_05_natural.png} \\
\multicolumn{7}{l}{\footnotesize\textit{Prompt (human\_05):} a young man with blonde spiky hair} \\
\addlinespace[2pt]
\includegraphics[width=0.135\textwidth]{figs/janus_comparison/no_adapter/human_06_natural.png} &
\includegraphics[width=0.135\textwidth]{figs/janus_comparison/adapter_v2/human_06_natural.png} &
\includegraphics[width=0.135\textwidth]{figs/janus_comparison/adapter_v3/human_06_natural.png} &
\includegraphics[width=0.135\textwidth]{figs/janus_comparison/adapter_v4/human_06_natural.png} &
\includegraphics[width=0.135\textwidth]{figs/janus_comparison/adapter_solo_gen/human_06_natural.png} &
\includegraphics[width=0.135\textwidth]{figs/janus_comparison/adapter_text_matched/human_06_natural.png} &
\includegraphics[width=0.135\textwidth]{figs/janus_comparison/adapter_text_far/human_06_natural.png} \\
\multicolumn{7}{l}{\footnotesize\textit{Prompt (human\_06):} a samurai warrior holding a katana} \\
\addlinespace[2pt]
\includegraphics[width=0.135\textwidth]{figs/janus_comparison/no_adapter/human_07_natural.png} &
\includegraphics[width=0.135\textwidth]{figs/janus_comparison/adapter_v2/human_07_natural.png} &
\includegraphics[width=0.135\textwidth]{figs/janus_comparison/adapter_v3/human_07_natural.png} &
\includegraphics[width=0.135\textwidth]{figs/janus_comparison/adapter_v4/human_07_natural.png} &
\includegraphics[width=0.135\textwidth]{figs/janus_comparison/adapter_solo_gen/human_07_natural.png} &
\includegraphics[width=0.135\textwidth]{figs/janus_comparison/adapter_text_matched/human_07_natural.png} &
\includegraphics[width=0.135\textwidth]{figs/janus_comparison/adapter_text_far/human_07_natural.png} \\
\multicolumn{7}{l}{\footnotesize\textit{Prompt (human\_07):} an old wizard with a long white beard} \\
\addlinespace[2pt]
\includegraphics[width=0.135\textwidth]{figs/janus_comparison/no_adapter/human_08_natural.png} &
\includegraphics[width=0.135\textwidth]{figs/janus_comparison/adapter_v2/human_08_natural.png} &
\includegraphics[width=0.135\textwidth]{figs/janus_comparison/adapter_v3/human_08_natural.png} &
\includegraphics[width=0.135\textwidth]{figs/janus_comparison/adapter_v4/human_08_natural.png} &
\includegraphics[width=0.135\textwidth]{figs/janus_comparison/adapter_solo_gen/human_08_natural.png} &
\includegraphics[width=0.135\textwidth]{figs/janus_comparison/adapter_text_matched/human_08_natural.png} &
\includegraphics[width=0.135\textwidth]{figs/janus_comparison/adapter_text_far/human_08_natural.png} \\
\multicolumn{7}{l}{\footnotesize\textit{Prompt (human\_08):} a teenage girl with pink hair} \\
\end{tabular}
\caption{Qualitative comparison across the seven training conditions on
Janus-Pro-7B. Columns 1--5 trace the image-only LoRA progression to the
\textit{Final} configuration; columns 6--7 add a joint text-training stream from
two semantically distant Wikipedia domains (Naruto franchise vs.\
physics/mathematics). All inferences share the same seed schedule; the adapter is
the only source of variation within a row.}
\label{fig:janus_qualitative_comparison}
\end{figure*}

\paragraph{Joint training: a corpus-size effect on perplexity.}
Adding a text stream inflates held-out text perplexity only when the text corpus
is small (Table~\ref{tab:joint}). The controlled far@320 run (distant domain,
matched chunk budget) reaches $55.1$, higher than the near-domain matched run,
while MMLU and POPE stay flat. The inflation is therefore a corpus-size effect (a
small corpus is revisited and over-fit), not a domain-distance effect, and it
does not correspond to any degradation in the other capability directions. This
matters for our reading of retention: a naive protocol could mistake this
perplexity inflation for forgetting.

\begin{table}[t]
\centering\small
\setlength{\tabcolsep}{6pt}
\begin{tabular}{llcccc}
\toprule
\textbf{Condition} & \textbf{Text corpus} & \textbf{PPL}$\downarrow$ & \textbf{MMLU}$\uparrow$ & \textbf{POPE acc}$\uparrow$ & \textbf{MME}$\uparrow$ \\
\midrule
Base (no adapter) & ---                  & $13.71$ & $0.492$ & $0.877$ & $1785$ \\
Final (image-only)& ---                  & $13.80$ & $0.491$ & $0.879$ & $1789$ \\
Joint matched     & near, 320 chunks     & $47.87$ & $0.505$ & $0.876$ & $1747$ \\
Joint far         & distant, 7672 chunks & $12.46$ & $0.491$ & $0.878$ & $1795$ \\
Joint far@320     & distant, 320 chunks  & $55.10$ & $0.481$ & $0.873$ & $1736$ \\
\bottomrule
\end{tabular}
\caption{Corpus-size control. Perplexity inflation tracks corpus size, not domain
distance; knowledge and understanding are unaffected.}
\label{tab:joint}
\end{table}


\section{Generalisation to a second visual style}
\label{app:sketch}
To test that the generation-direction result is not specific to one style, we
replicate the image protocol on a second, unrelated dataset: black-and-white
line-sketch scenes (\texttt{zoheb/sketch-scene}), using the \emph{same}
configuration as Naruto ($r{=}48$, $\alpha{=}96$, $\lambda{=}0.5$, LR
$10^{-4}$). The core result holds (Table~\ref{tab:sketch}): the adapter acquires
the style at the distribution level (CMMD $0.253\!\to\!0.159$, KID
$0.263\!\to\!0.096$, DINO similarity to the sketch reference
$0.133\!\to\!0.371$), while the language and knowledge probes show no
catastrophic forgetting (WikiText PPL $13.71\!\to\!14.18$, MMLU
$0.492\!\to\!0.479$).

This replication also surfaces two findings we use in the main text. First,
\emph{the PEFT configuration is dataset-dependent}: the recipe that saturates a
high-frequency, colour-saturated style only partially imprints the sparse,
low-frequency sketch style (Figure~\ref{fig:sketch}). The per-prompt edit
strength is about half (CLIP DirSim $0.086$, versus $\approx\!0.15$ for Naruto at
the same $\lambda$), and it carries a measurable image-quality cost (CLIP-IQA
$0.747\!\to\!0.592$). Scaling the recipe does not help: a stronger configuration
(rank $64$, $\lambda{=}0.25$) moves the aggregate distribution closer to the
reference (KID $0.096\!\to\!0.065$, DINO $0.371\!\to\!0.521$) but does not
improve the per-prompt edit strength (DirSim $0.086\!\to\!0.089$) and lowers
image quality (CLIP-IQA $0.59\!\to\!0.49$), indicating the limit is intrinsic to
the low-frequency style rather than a matter of capacity. Second, this is the
target whose retention footprint is uniformly the largest of all five, and most
visible on MME ($1785\!\to\!1713$, versus flat for Naruto): it is the third of
the three convergent settings discussed in Section~\ref{sec:boundary}.

\begin{table}[t]
\centering\small
\setlength{\tabcolsep}{8pt}
\begin{tabular}{lcc}
\toprule
Metric & Base & Sketch LoRA \\
\midrule
\multicolumn{3}{l}{\textit{Acquisition (vs.\ sketch reference, $n{=}80$)}}\\
CMMD $\downarrow$      & $.253$ & $\mathbf{.159}_{\pm.007}$ \\
KID $\downarrow$       & $.263$ & $\mathbf{.096}_{\pm.013}$ \\
DINO sim $\uparrow$    & $.133$ & $\mathbf{.371}_{\pm.021}$ \\
CLIP DirSim $\uparrow$ & ---    & $.086_{\pm.006}$ \\
\midrule
\multicolumn{3}{l}{\textit{Retention}}\\
WikiText PPL $\downarrow$ & $13.71_{\pm.21}$ & $14.18_{\pm.22}$ \\
MMLU $\uparrow$           & $.492_{\pm.013}$ & $.479_{\pm.013}$ \\
POPE acc $\uparrow$       & $.877_{\pm.006}$ & $.871_{\pm.006}$ \\
MME $\uparrow$            & $1785_{\pm31}$   & $1713_{\pm31}$ \\
\bottomrule
\end{tabular}
\caption{Generalisation to a second image dataset (line-sketch scenes,
\texttt{zoheb/sketch-scene}) with the \emph{same} configuration as Naruto. The
style is acquired at the distribution level and the language and knowledge probes
show no catastrophic forgetting; the per-prompt edit strength (CLIP DirSim
$0.086$) is about half of Naruto's, evidence that the PEFT recipe is
dataset-dependent. This is also the largest retention footprint of the five
targets, most visible on MME. FID/ArtFID are omitted (numerically unstable at
$n{=}80$).}
\label{tab:sketch}
\end{table}

\begin{figure}[t]
\centering
\begin{subfigure}{0.24\textwidth}
\includegraphics[width=\linewidth]{figs/sketch_comparison/zebra_before.png}
\caption*{\footnotesize Base}
\end{subfigure}
\begin{subfigure}{0.24\textwidth}
\includegraphics[width=\linewidth]{figs/sketch_comparison/zebra_after.png}
\caption*{\footnotesize Sketch LoRA}
\end{subfigure}
\begin{subfigure}{0.24\textwidth}
\includegraphics[width=\linewidth]{figs/sketch_comparison/kite_before.png}
\caption*{\footnotesize Base}
\end{subfigure}
\begin{subfigure}{0.24\textwidth}
\includegraphics[width=\linewidth]{figs/sketch_comparison/kite_after.png}
\caption*{\footnotesize Sketch LoRA}
\end{subfigure}
\caption{Generalisation to line-sketch scenes (same prompt and seed, base vs.\
adapter). The style transfers for some prompts (left pair: white canvas, ink
linework) but only partially for others (right pair: sketch strokes appear yet
photographic colour persists), consistent with the weaker, dataset-dependent
acquisition in Table~\ref{tab:sketch}.}
\label{fig:sketch}
\end{figure}


\section{Qualitative grids}
\label{app:grids}

\begin{figure*}[p]
\centering
\setlength{\tabcolsep}{1.5pt}\small
\renewcommand{\arraystretch}{1.0}
\begin{tabular}{ccccccc}
\textbf{Base} & \textbf{r=2} & \textbf{r=4} & \textbf{r=8} & \textbf{r=16} & \textbf{r=32} & \textbf{r=64} \\
\midrule
\includegraphics[width=0.135\textwidth]{figs/rank_comparison/base/eval_00_natural.png} &
\includegraphics[width=0.135\textwidth]{figs/rank_comparison/r2/eval_00_natural.png} &
\includegraphics[width=0.135\textwidth]{figs/rank_comparison/r4/eval_00_natural.png} &
\includegraphics[width=0.135\textwidth]{figs/rank_comparison/r8/eval_00_natural.png} &
\includegraphics[width=0.135\textwidth]{figs/rank_comparison/r16/eval_00_natural.png} &
\includegraphics[width=0.135\textwidth]{figs/rank_comparison/r32/eval_00_natural.png} &
\includegraphics[width=0.135\textwidth]{figs/rank_comparison/r64/eval_00_natural.png} \\
\multicolumn{7}{l}{\footnotesize\textit{Prompt:} a red panda with a black face and brown eyes} \\
\addlinespace[2pt]
\includegraphics[width=0.135\textwidth]{figs/rank_comparison/base/eval_01_natural.png} &
\includegraphics[width=0.135\textwidth]{figs/rank_comparison/r2/eval_01_natural.png} &
\includegraphics[width=0.135\textwidth]{figs/rank_comparison/r4/eval_01_natural.png} &
\includegraphics[width=0.135\textwidth]{figs/rank_comparison/r8/eval_01_natural.png} &
\includegraphics[width=0.135\textwidth]{figs/rank_comparison/r16/eval_01_natural.png} &
\includegraphics[width=0.135\textwidth]{figs/rank_comparison/r32/eval_01_natural.png} &
\includegraphics[width=0.135\textwidth]{figs/rank_comparison/r64/eval_01_natural.png} \\
\multicolumn{7}{l}{\footnotesize\textit{Prompt:} a man with glasses and a hoodie} \\
\addlinespace[2pt]
\includegraphics[width=0.135\textwidth]{figs/rank_comparison/base/human_05_natural.png} &
\includegraphics[width=0.135\textwidth]{figs/rank_comparison/r2/human_05_natural.png} &
\includegraphics[width=0.135\textwidth]{figs/rank_comparison/r4/human_05_natural.png} &
\includegraphics[width=0.135\textwidth]{figs/rank_comparison/r8/human_05_natural.png} &
\includegraphics[width=0.135\textwidth]{figs/rank_comparison/r16/human_05_natural.png} &
\includegraphics[width=0.135\textwidth]{figs/rank_comparison/r32/human_05_natural.png} &
\includegraphics[width=0.135\textwidth]{figs/rank_comparison/r64/human_05_natural.png} \\
\multicolumn{7}{l}{\footnotesize\textit{Prompt:} a young man with blonde spiky hair} \\
\addlinespace[2pt]
\includegraphics[width=0.135\textwidth]{figs/rank_comparison/base/human_07_natural.png} &
\includegraphics[width=0.135\textwidth]{figs/rank_comparison/r2/human_07_natural.png} &
\includegraphics[width=0.135\textwidth]{figs/rank_comparison/r4/human_07_natural.png} &
\includegraphics[width=0.135\textwidth]{figs/rank_comparison/r8/human_07_natural.png} &
\includegraphics[width=0.135\textwidth]{figs/rank_comparison/r16/human_07_natural.png} &
\includegraphics[width=0.135\textwidth]{figs/rank_comparison/r32/human_07_natural.png} &
\includegraphics[width=0.135\textwidth]{figs/rank_comparison/r64/human_07_natural.png} \\
\multicolumn{7}{l}{\footnotesize\textit{Prompt:} an old wizard with a long white beard} \\
\addlinespace[2pt]
\includegraphics[width=0.135\textwidth]{figs/rank_comparison/base/human_08_natural.png} &
\includegraphics[width=0.135\textwidth]{figs/rank_comparison/r2/human_08_natural.png} &
\includegraphics[width=0.135\textwidth]{figs/rank_comparison/r4/human_08_natural.png} &
\includegraphics[width=0.135\textwidth]{figs/rank_comparison/r8/human_08_natural.png} &
\includegraphics[width=0.135\textwidth]{figs/rank_comparison/r16/human_08_natural.png} &
\includegraphics[width=0.135\textwidth]{figs/rank_comparison/r32/human_08_natural.png} &
\includegraphics[width=0.135\textwidth]{figs/rank_comparison/r64/human_08_natural.png} \\
\multicolumn{7}{l}{\footnotesize\textit{Prompt:} a teenage girl with pink hair} \\
\end{tabular}
\caption{Qualitative rank sweep (Final configuration): the target style is already
present at $r{=}2$ and only refines with rank, remaining stable across the range.}
\label{fig:rank_visual}
\end{figure*}

\begin{figure}[t]
\centering
\includegraphics[width=\textwidth]{figs/rank_sweep/rank_sweep_figure.pdf}
\caption{Rank sweep on the Final configuration (absolute metric values).
\textbf{Left:} visual-style acquisition (CLIP directional similarity); the
dotted line at $0$ marks an adapter with no directional effect. Acquisition is
positive at every rank and rises modestly, saturating around $r=16$--$32$.
\textbf{Right:} visual understanding (POPE accuracy); the dashed line is the
base model. The adapter tracks the base across the whole range, i.e.\ no
cross-modal interference at any rank.}
\label{fig:rank_sweep_panels}
\end{figure}

\begin{figure}[t]
\centering
\includegraphics[width=0.7\columnwidth]{figs/kl_sweep/kl_sweep_image.pdf}
\caption{Style acquisition (CLIP directional similarity, absolute scale)
versus the KL weight $\lambda$ at rank $32$. The dotted line at $0$ marks an
adapter with no directional effect. Acquisition peaks at
$\lambda\!\approx\!0.25$--$0.5$ and collapses toward $0$ for $\lambda\!\geq\!1$,
where the adapter is over-regularised toward the base.}
\label{fig:kl_curve}
\end{figure}

\begin{figure*}[p]
\centering
\setlength{\tabcolsep}{1.5pt}\small
\begin{tabular}{cccccc}
\textbf{Base} & \textbf{$\lambda{=}0$} & \textbf{$\lambda{=}0.25$} & \textbf{$\lambda{=}0.5$} & \textbf{$\lambda{=}1.0$} & \textbf{$\lambda{=}1.5$} \\
\midrule
\includegraphics[width=0.155\textwidth]{figs/kl_comparison/base/eval_00_natural.png} &
\includegraphics[width=0.155\textwidth]{figs/kl_comparison/kl0/eval_00_natural.png} &
\includegraphics[width=0.155\textwidth]{figs/kl_comparison/kl025/eval_00_natural.png} &
\includegraphics[width=0.155\textwidth]{figs/kl_comparison/kl05/eval_00_natural.png} &
\includegraphics[width=0.155\textwidth]{figs/kl_comparison/kl1/eval_00_natural.png} &
\includegraphics[width=0.155\textwidth]{figs/kl_comparison/kl15/eval_00_natural.png} \\
\multicolumn{6}{l}{\footnotesize\textit{Prompt:} a red panda with a black face and brown eyes} \\
\addlinespace[2pt]
\includegraphics[width=0.155\textwidth]{figs/kl_comparison/base/eval_01_natural.png} &
\includegraphics[width=0.155\textwidth]{figs/kl_comparison/kl0/eval_01_natural.png} &
\includegraphics[width=0.155\textwidth]{figs/kl_comparison/kl025/eval_01_natural.png} &
\includegraphics[width=0.155\textwidth]{figs/kl_comparison/kl05/eval_01_natural.png} &
\includegraphics[width=0.155\textwidth]{figs/kl_comparison/kl1/eval_01_natural.png} &
\includegraphics[width=0.155\textwidth]{figs/kl_comparison/kl15/eval_01_natural.png} \\
\multicolumn{6}{l}{\footnotesize\textit{Prompt:} a man with glasses and a hoodie} \\
\addlinespace[2pt]
\includegraphics[width=0.155\textwidth]{figs/kl_comparison/base/human_05_natural.png} &
\includegraphics[width=0.155\textwidth]{figs/kl_comparison/kl0/human_05_natural.png} &
\includegraphics[width=0.155\textwidth]{figs/kl_comparison/kl025/human_05_natural.png} &
\includegraphics[width=0.155\textwidth]{figs/kl_comparison/kl05/human_05_natural.png} &
\includegraphics[width=0.155\textwidth]{figs/kl_comparison/kl1/human_05_natural.png} &
\includegraphics[width=0.155\textwidth]{figs/kl_comparison/kl15/human_05_natural.png} \\
\multicolumn{6}{l}{\footnotesize\textit{Prompt:} a young man with blonde spiky hair} \\
\addlinespace[2pt]
\includegraphics[width=0.155\textwidth]{figs/kl_comparison/base/human_07_natural.png} &
\includegraphics[width=0.155\textwidth]{figs/kl_comparison/kl0/human_07_natural.png} &
\includegraphics[width=0.155\textwidth]{figs/kl_comparison/kl025/human_07_natural.png} &
\includegraphics[width=0.155\textwidth]{figs/kl_comparison/kl05/human_07_natural.png} &
\includegraphics[width=0.155\textwidth]{figs/kl_comparison/kl1/human_07_natural.png} &
\includegraphics[width=0.155\textwidth]{figs/kl_comparison/kl15/human_07_natural.png} \\
\multicolumn{6}{l}{\footnotesize\textit{Prompt:} an old wizard with a long white beard} \\
\addlinespace[2pt]
\includegraphics[width=0.155\textwidth]{figs/kl_comparison/base/human_08_natural.png} &
\includegraphics[width=0.155\textwidth]{figs/kl_comparison/kl0/human_08_natural.png} &
\includegraphics[width=0.155\textwidth]{figs/kl_comparison/kl025/human_08_natural.png} &
\includegraphics[width=0.155\textwidth]{figs/kl_comparison/kl05/human_08_natural.png} &
\includegraphics[width=0.155\textwidth]{figs/kl_comparison/kl1/human_08_natural.png} &
\includegraphics[width=0.155\textwidth]{figs/kl_comparison/kl15/human_08_natural.png} \\
\multicolumn{6}{l}{\footnotesize\textit{Prompt:} a teenage girl with pink hair} \\
\end{tabular}
\caption{Effect of the KL-to-base weight $\lambda$ on visual style acquisition
(rank $32$). Style transfer is strong for $\lambda\!\in\![0,0.5]$ and visibly
collapses back toward the base for $\lambda\!\geq\!1$, where the adapter is
over-regularised.}
\label{fig:kl_visual}
\end{figure*}


\section{Text-direction rank sweep}
\label{app:textsweep}
Table~\ref{tab:textrank} sweeps the LoRA rank on the poetry target at
$\lambda{=}0.5$, the plasticity axis complementary to the anchor sweep of
Table~\ref{tab:textkl}. Increasing the rank from $4$ to $64$ steadily improves
acquisition (poetry perplexity $32.7$ to $23.8$) while MMLU, POPE and MME stay at
the base level at every rank, and WikiText perplexity is not degraded and is in
fact slightly below the base throughout. Specializing the text direction
therefore leaves the other capabilities intact across capacity, mirroring the
image result of Table~\ref{tab:rank}. Pushing the rank further to $128$ confirms
diminishing returns: acquisition improves only marginally (poetry perplexity
$23.8$ to $23.2$), while a first mild understanding cost appears (MME $1122$ to
$1105$, POPE $0.877$ to $0.867$); this is the second of the three convergent
settings of Section~\ref{sec:boundary}. We therefore report up to rank $64$,
which sits at the knee of this trade-off.

\begin{table}[t]
\centering\small
\setlength{\tabcolsep}{7pt}
\begin{tabular}{lcccccc}
\toprule
& Base & $r{=}4$ & $r{=}8$ & $r{=}16$ & $r{=}32$ & $r{=}64$ \\
\midrule
\multicolumn{7}{l}{\textit{Acquisition (target style)}}\\
Poetry PPL $\downarrow$ & $32.7$ & $25.6$ & $25.3$ & $24.9$ & $24.4$ & $\mathbf{23.8}$ \\
\midrule
\multicolumn{7}{l}{\textit{Retention}}\\
WikiText PPL $\downarrow$ & $13.71$ & $12.52$ & $\underline{12.51}$ & $\underline{12.51}$ & $12.53$ & $12.59$ \\
MMLU $\uparrow$           & $.481$ & $.481$ & $.477$ & $.479$ & $.477$ & $.488$ \\
POPE acc $\uparrow$       & $.875$ & $.871$ & $.871$ & $.871$ & $.873$ & $.877$ \\
MME $\uparrow$            & $1130$ & $1129$ & $1111$ & $1128$ & $1130$ & $1122$ \\
\bottomrule
\end{tabular}
\caption{Text (poetry) rank sweep at $\lambda{=}0.5$. The style is acquired (bold
= best); knowledge and understanding stay flat. WikiText perplexity is not
degraded and is in fact slightly lower at every rank (underlined best).
Subsampled probe configuration, single seed.}
\label{tab:textrank}
\end{table}


\section{Image generation under the text adapter}
\label{app:t2itext}
\begin{figure*}[t]
\centering
\setlength{\tabcolsep}{1.5pt}
\renewcommand{\arraystretch}{0.9}
\small
\begin{tabular}{cccccc}
\textbf{Base} & $\lambda{=}0$ & $\lambda{=}0.25$ & $\lambda{=}0.5$ & $\lambda{=}1$ & $\lambda{=}1.5$ \\
\midrule
\includegraphics[width=0.15\textwidth]{figs/t2i_text/kl0/prompt0_base.png} &
\includegraphics[width=0.15\textwidth]{figs/t2i_text/kl0/prompt0_adapter.png} &
\includegraphics[width=0.15\textwidth]{figs/t2i_text/kl025/prompt0_adapter.png} &
\includegraphics[width=0.15\textwidth]{figs/t2i_text/kl05/prompt0_adapter.png} &
\includegraphics[width=0.15\textwidth]{figs/t2i_text/kl1/prompt0_adapter.png} &
\includegraphics[width=0.15\textwidth]{figs/t2i_text/kl15/prompt0_adapter.png} \\
\multicolumn{6}{l}{\footnotesize\textit{Prompt:} a cat sitting on a chair} \\
\addlinespace[2pt]
\includegraphics[width=0.15\textwidth]{figs/t2i_text/kl0/prompt1_base.png} &
\includegraphics[width=0.15\textwidth]{figs/t2i_text/kl0/prompt1_adapter.png} &
\includegraphics[width=0.15\textwidth]{figs/t2i_text/kl025/prompt1_adapter.png} &
\includegraphics[width=0.15\textwidth]{figs/t2i_text/kl05/prompt1_adapter.png} &
\includegraphics[width=0.15\textwidth]{figs/t2i_text/kl1/prompt1_adapter.png} &
\includegraphics[width=0.15\textwidth]{figs/t2i_text/kl15/prompt1_adapter.png} \\
\multicolumn{6}{l}{\footnotesize\textit{Prompt:} a medieval castle on a hill} \\
\addlinespace[2pt]
\includegraphics[width=0.15\textwidth]{figs/t2i_text/kl0/prompt2_base.png} &
\includegraphics[width=0.15\textwidth]{figs/t2i_text/kl0/prompt2_adapter.png} &
\includegraphics[width=0.15\textwidth]{figs/t2i_text/kl025/prompt2_adapter.png} &
\includegraphics[width=0.15\textwidth]{figs/t2i_text/kl05/prompt2_adapter.png} &
\includegraphics[width=0.15\textwidth]{figs/t2i_text/kl1/prompt2_adapter.png} &
\includegraphics[width=0.15\textwidth]{figs/t2i_text/kl15/prompt2_adapter.png} \\
\multicolumn{6}{l}{\footnotesize\textit{Prompt:} a bowl of fruit on a table} \\
\addlinespace[2pt]
\includegraphics[width=0.15\textwidth]{figs/t2i_text/kl0/prompt3_base.png} &
\includegraphics[width=0.15\textwidth]{figs/t2i_text/kl0/prompt3_adapter.png} &
\includegraphics[width=0.15\textwidth]{figs/t2i_text/kl025/prompt3_adapter.png} &
\includegraphics[width=0.15\textwidth]{figs/t2i_text/kl05/prompt3_adapter.png} &
\includegraphics[width=0.15\textwidth]{figs/t2i_text/kl1/prompt3_adapter.png} &
\includegraphics[width=0.15\textwidth]{figs/t2i_text/kl15/prompt3_adapter.png} \\
\multicolumn{6}{l}{\footnotesize\textit{Prompt:} an astronaut riding a horse} \\
\end{tabular}
\caption{Image generation under the text (poetry) adapter across KL weights
(rank 16). Column 1 is the base model; columns 2 to 6 are the poetry adapter at
increasing $\lambda$, on the same prompts and seeds. The generated images stay
close to the base across the whole range, the qualitative counterpart of the
near-zero Gen-CMMD in Table~\ref{tab:textkl}.}
\label{fig:t2i_text_kl}
\end{figure*}


\section{Intervention strength and its weight-space localization}
\label{app:dw}
The aggregate intervention strengths are reported in Table~\ref{tab:dw} of the
main text. The base norm $\lVert W\rVert_F$ is read from the frozen weights;
$\lVert\Delta W\rVert_F=\lVert(\alpha/r)BA\rVert_F$ is computed from the LoRA
factors via the $r\times r$ trace identity, so the full $\Delta W$ is never
formed. Figure~\ref{fig:dw-heatmap} resolves the same quantity per module, which
is what makes the ``localized'' claim concrete rather than aggregate.

\begin{figure*}[t]
\centering
\begin{subfigure}{0.32\linewidth}\centering
  \includegraphics[width=\linewidth]{figs/weight_delta/dw_heatmap_img_r16.png}
  \caption{Image gen., $r{=}16$}
\end{subfigure}\hfill
\begin{subfigure}{0.32\linewidth}\centering
  \includegraphics[width=\linewidth]{figs/weight_delta/dw_heatmap_text_r16.png}
  \caption{Text/poetry, $r{=}16$}
\end{subfigure}\hfill
\begin{subfigure}{0.32\linewidth}\centering
  \includegraphics[width=\linewidth]{figs/weight_delta/dw_heatmap_text_r64.png}
  \caption{Text/poetry, $r{=}64$}
\end{subfigure}
\caption{Per-module relative weight change
$\lVert\Delta W\rVert_F/\lVert W\rVert_F$ across layers (rows) and the seven
projections (columns). Common to both directions, the update spares the
attention query/key/value projections almost everywhere, awakening only in the
last few layers, and concentrates in the MLP. The directions differ in the fine
locus: text spreads the change broadly over \texttt{gate} across all depths,
whereas image concentrates it toward the output, peaking at the final layer's
\texttt{down} projection. Raising the rank (text $r{=}16$ vs.\ $r{=}64$) scales
the same pattern up rather than relocating it.}
\label{fig:dw-heatmap}
\end{figure*}


\section{Regularizer ablation}
\label{app:reg}
We compare the KL-to-base regularizer at fixed rank 32 in the image-generation
direction against three alternatives: L2-SP \citep{xuhong2018l2sp} (an L2 penalty
on the LoRA update, weight space), logit-MSE (a squared-error match to the base
logits, output space), and the reconstruction-loss objective of Anole
\citep{chern2024anole}, the standard recipe for fine-tuning native generators
(\texttt{CE}$+0.5\,$recon, no base anchor). All share the same training budget and
data and are evaluated identically on $80$ held-out prompts under the canonical,
trigger-free protocol used in the rest of the paper; we also report each adapter's
intervention strength $\lVert\Delta W\rVert_F/\lVert W\rVert_F$. The KL attains the
best distributional acquisition of all methods: at essentially the same
intervention strength as the unregularized baseline ($1.57$ vs.\ $1.55\%$) it
reaches a markedly lower CMMD ($.108$ vs.\ $.155$), i.e.\ it \emph{redirects} the
update rather than merely shrinking it. Logit-MSE over-constrains the output
distribution: it keeps image quality high but barely transfers the style without a
prompt trigger (DirSim $.029$, the worst of all methods), which the prompt-triggered
evaluation had masked. Weight-space L2-SP simply shrinks the update (to $0.60\%$).
The Anole reconstruction loss acquires the style comparably to the unregularized
baseline (DirSim $.162$) but with worse distributional fidelity than the KL (CMMD
$.148$ vs.\ $.108$) and, being an acquisition term rather than a base anchor, offers
no forgetting protection. The KL alone attains the best acquisition while keeping
the adapter anchored to the base. Single-seed run.

\begin{table}[t]\centering\small
\setlength{\tabcolsep}{3.5pt}
\begin{tabular}{lcccccc}
\toprule
Regularizer & $\lVert\Delta W\rVert/\lVert W\rVert$ & CMMD $\downarrow$ & KID $\downarrow$ & DINO $\uparrow$ & CLIP-IQA $\uparrow$ & DirSim $\uparrow$ \\
\midrule
None ($\lambda{=}0$) & $1.55\%$ & $.155_{\pm.009}$ & $.027_{\pm.004}$ & $.626_{\pm.013}$ & $.666_{\pm.011}$ & $.157_{\pm.007}$ \\
L2-SP                & $0.60\%$ & $.139_{\pm.009}$ & $.018_{\pm.004}$ & $.624_{\pm.011}$ & $.683_{\pm.008}$ & $\mathbf{.164}_{\pm.007}$ \\
Logit-MSE            & $1.14\%$ & $.216_{\pm.012}$ & $.106_{\pm.011}$ & $.410_{\pm.018}$ & $\mathbf{.725}_{\pm.014}$ & $.029_{\pm.006}$ \\
KL-to-base (ours)    & $1.57\%$ & $\mathbf{.108}_{\pm.010}$ & $\mathbf{.014}_{\pm.004}$ & $\mathbf{.643}_{\pm.013}$ & $.691_{\pm.010}$ & $.153_{\pm.007}$ \\
Anole (recon)       & $1.56\%$ & $.148_{\pm.009}$ & $.026_{\pm.005}$ & $.637_{\pm.011}$ & $.666_{\pm.010}$ & $.162_{\pm.007}$ \\
\bottomrule
\end{tabular}
\caption{Regularizer ablation at rank 32 (image direction, $N{=}80$, single seed;
$\pm$ one standard error; bold = best per metric). Under the canonical,
trigger-free evaluation, the KL attains the best distributional acquisition
(CMMD, KID, DINO), separating clearly from the no-regularizer baseline at matched
intervention strength ($\lVert\Delta W\rVert/\lVert W\rVert$ $1.57$ vs.\ $1.55\%$,
CMMD $.108$ vs.\ $.155$). Logit-MSE over-constrains the output distribution and
barely transfers the style without the prompt trigger (DirSim $.029$) despite high
image quality; L2-SP merely shrinks the update. The Anole reconstruction loss
acquires the style comparably to the unregularized baseline but with worse
distributional fidelity than the KL and no base anchoring.}
\label{tab:reg-ablation}
\end{table}

\clearpage
\section{Metric debugging: is the $\lambda$ optimum a bug, memorization, or noise?}
\label{app:metricdebug}
The style metrics put the image-generation acquisition optimum at $\lambda{=}0.25$,
while the eye mildly prefers $\lambda{=}0.5$. We resolve this directly by
re-scoring the \emph{already-generated} KL-sweep images (no retraining) against
both held-out and training references, and with two additional embedding metrics
(Table~\ref{tab:metric-debug}). \textbf{Not noise:} the $\lambda$-ranking of CMMD
is identical against held-out and training reals (both order
$0.25\!<\!0.5\!<\!0\!<\!1\!<\!1.5$), so the optimum is stable, not a small-$N$
artefact. \textbf{Not memorization or leakage:} for every $\lambda$ the
nearest-neighbour DINOv2 similarity to \emph{training} images is \emph{lower} than
to held-out images (and CMMD to training is higher), so the generations are not
near-copies of the training set. \textbf{A metric-choice effect, not a bug:} CMMD
(a CLIP-space distribution distance) favours $\lambda{=}0.25$, but the
memorization-robust DINOv2-space distribution distance (DINO-MMD) instead favours
$\lambda{=}0.5$ ($.035$ vs.\ $.041$), and content preservation -- the DINOv2
similarity between the adapter's image and the base image for the same prompt --
rises from $.32$ at $\lambda{=}0.25$ to $.39$ at $\lambda{=}0.5$. The CLIP style
metrics reward pure style proximity (optimum $\lambda{=}0.25$); the DINO
distribution distance and content preservation reward style transfer that keeps
the prompt content (optimum $\lambda{=}0.5$), which is what the eye weighs, and
what an independent VLM judge confirms (Table~\ref{tab:judge}). This justifies the
$\lambda{=}0.5$ Final recipe.

\begin{table}[h]
\centering\small
\setlength{\tabcolsep}{5pt}
\begin{tabular}{lcccccc}
\toprule
$\lambda$ & CMMD$_{\text{test}}\downarrow$ & CMMD$_{\text{train}}\downarrow$ & DINO-MMD$\downarrow$ & NN$_{\text{test}}\uparrow$ & NN$_{\text{train}}\uparrow$ & CtntPres$\uparrow$ \\
\midrule
$0$    & $.155$ & $.163$ & $.044$ & $.626$ & $.609$ & $.303$ \\
$0.25$ & $\mathbf{.100}$ & $.111$ & $.041$ & $.647$ & $.634$ & $.317$ \\
$0.5$  & $.108$ & $.124$ & $\mathbf{.035}$ & $.643$ & $.633$ & $.389$ \\
$1$    & $.205$ & $.221$ & $.091$ & $.439$ & $.411$ & $.721$ \\
$1.5$  & $.221$ & $.235$ & $.114$ & $.388$ & $.364$ & $.792$ \\
\bottomrule
\end{tabular}
\caption{Metric debugging on the KL sweep (rank 32), re-scoring the
already-generated images (no retraining). CMMD$_{\text{test/train}}$: CLIP-space
distribution distance to held-out vs.\ training reals (identical ranking $\Rightarrow$
signal not noise; train $>$ test $\Rightarrow$ no training over-fit). NN: nearest
DINOv2 similarity to held-out vs.\ training reals (train $<$ test everywhere
$\Rightarrow$ no memorization). DINO-MMD: memorization-robust DINOv2-space
distribution distance -- optimum $\lambda{=}0.5$. CtntPres: DINOv2 similarity
between the adapter and base image for the same prompt -- rises with $\lambda$. The
CLIP metrics favour $\lambda{=}0.25$; DINO-MMD and content preservation favour
$\lambda{=}0.5$.}
\label{tab:metric-debug}
\end{table}

As an independent corroboration we add an open VLM judge (Qwen2-VL-7B, not Janus,
to avoid self-evaluation), comparing $\lambda{=}0.25$ vs.\ $\lambda{=}0.5$ pairwise
on the same $80$ prompts in both image orders and counting only order-consistent
preferences (MT-Bench style). It prefers $\lambda{=}0.5$ on every criterion,
significantly on content (Table~\ref{tab:judge}).

\begin{table}[t]
\centering\small
\setlength{\tabcolsep}{8pt}
\begin{tabular}{lccc}
\toprule
Criterion & $\lambda{=}0.25$ win & $\lambda{=}0.5$ win & tie \\
\midrule
Style    & $23.8\%$ & $\mathbf{30.0\%}$ & $46.2\%$ \\
Content  & $6.2\%$  & $\mathbf{20.0\%}$ & $73.8\%$ \\
Overall  & $23.8\%$ & $\mathbf{30.0\%}$ & $46.2\%$ \\
\bottomrule
\end{tabular}
\caption{Independent VLM judge (Qwen2-VL-7B), $\lambda{=}0.25$ vs.\ $\lambda{=}0.5$
on $80$ held-out prompts, each scored in both image orders. A configuration wins a
prompt only if preferred in \emph{both} orders; order-dependent results are ties
(position-bias mitigation). On the order-consistent comparisons the judge prefers
$\lambda{=}0.5$ on every criterion, significantly on content (decisive win-rate
$76\%$, $95\%$ CI $[55,89]$), and never favours $\lambda{=}0.25$.}
\label{tab:judge}
\end{table}

\clearpage
\section{Text-direction qualitative examples}
\label{app:textqual}
We show base-vs-adapter continuations for the poetry text-LoRA ($r{=}16$,
$\lambda{=}0.5$); the base is the identical model with the adapter disabled, so the
only change between columns is the LoRA. Decoding is greedy with a repetition
penalty of $1.3$ (Appendix~\ref{app:impl}). Each incipit is continued from a fixed
two-line verse couplet given as hidden context to \emph{both} models (so the
unified model stays in a poem rather than its assistant persona); we display only
the text generated after the incipit. With the same primed context the adapter
sustains metrical, lineated verse, whereas the base drifts into prose or stops
short. On bare, non-poetic prompts the register does not surface: the LoRA is a
next-token language model, not instruction-tuned.

\begin{table}[t]\centering\small
\renewcommand{\arraystretch}{1.3}
\begin{tabular}{@{}p{0.45\linewidth}p{0.45\linewidth}@{}}
\toprule
\multicolumn{2}{@{}l}{\textit{Incipit:} ``The autumn wind\,\dots''}\\
\midrule
\textbf{Base (adapter off)} & \textbf{Poetry LoRA}\\[2pt]
\footnotesize is blowing in our faces; it's a time to reflect on what has been
accomplished this year\,\dots
&
\footnotesize is blowing in our faces; it's a-blowing! ,  O what shall I do?\\
\midrule
\multicolumn{2}{@{}l}{\textit{Incipit:} ``Upon the silent hill there\,\dots''}\\
\midrule
\textbf{Base (adapter off)} & \textbf{Poetry LoRA}\\[2pt]
\footnotesize lies a light of gold; ,  O! how it gleams in twilight's shade: it
is not day nor night but something more divine.
&
\footnotesize stands a tree;\newline
The leaves are green as emeralds in their sleep:\newline
And on its branches hang an empty cage,\newline
A prison for some bird of sorrow's keep!\\
\midrule
\multicolumn{2}{@{}l}{\textit{Incipit:} ``Beneath the ancient oak,\,\dots''}\\
\midrule
\textbf{Base (adapter off)} & \textbf{Poetry LoRA}\\[2pt]
\footnotesize where time has worn its way; --- O how I love thee!
&
\footnotesize where I have stood;\newline
The sun was setting on a summer's day: ---\newline
I saw her face for one brief moment flash!\newline
And then she vanished from before mine eyes.\\
\midrule
\multicolumn{2}{@{}l}{\textit{Incipit:} ``O gentle heart, what\,\dots''}\\
\midrule
\textbf{Base (adapter off)} & \textbf{Poetry LoRA}\\[2pt]
\footnotesize words of love can say? --- O! for a thousand tongues to tell thee how
I feel this day; but none are fit nor true \dots\ (1 Samuel 25).
&
\footnotesize is thy secret? What art thou made of but a little child's delight in
play; \dots\ O say not ``I am nothing,'' for I see how much you have to give: your
eyes shine like stars on earth when they look at me\,\dots\\
\bottomrule
\end{tabular}
\caption{Text-direction qualitative examples (poetry LoRA, $r{=}16$,
$\lambda{=}0.5$; greedy decoding, repetition penalty $1.3$). Each incipit is
continued from a fixed verse couplet given as hidden context to \emph{both}
models; we show only the text after the incipit. With the same primed context the
adapter sustains lineated, imagistic verse (e.g.\ the emerald-leaf quatrain),
whereas the base is shorter, more abstract, or drifts into prose (top) or
scriptural citation (bottom).}
\label{tab:text-examples}
\end{table}


\section{Medicine: qualitative before and after}
\label{app:medicine}
Same prompts, same model, adapter off (base) vs.\ on (medicine LoRA). The text
continuations show the register shift toward clinical prose; the image grid shows
that text specialisation leaves the generation direction essentially unchanged.

\begin{table}[t]
\centering\footnotesize
\setlength{\tabcolsep}{6pt}\renewcommand{\arraystretch}{1.25}
\begin{tabular}{p{0.45\textwidth} p{0.45\textwidth}}
\toprule
\textbf{Base} & \textbf{Medicine LoRA} \\
\midrule
\multicolumn{2}{@{}p{0.95\textwidth}@{}}{\textit{Prompt:} ``What is the mechanism of action of metformin?''}\\
Metformin works by inhibiting glucose production in cells, particularly those
that are most active at high concentrations. Hello! I'm here to help you with
your question [\dots]
&
Metformin works by inhibiting glucose uptake in muscle cells and reducing
hepatic gluconeogenesis. [\dots] \\
\midrule
\multicolumn{2}{@{}p{0.95\textwidth}@{}}{\textit{Prompt:} ``A 58-year-old man with hypertension presents with acute chest pain radiating to the left arm; the most likely diagnosis is\dots''}\\
\textit{(returns another question instead of a diagnosis:)} ``What are some
common causes of sudden cardiac death in older adults?''
&
unstable angina. A coronary catheterization should be performed as soon as
possible and an appropriate treatment plan must begin immediately; if there are
no signs of myocardial infarction (MI), an electrocardiogram (ECG) and stress
test [\dots] \\
\midrule
\multicolumn{2}{@{}p{0.95\textwidth}@{}}{\textit{Prompt (general-knowledge control):} ``Who wrote the play Romeo and Juliet?''}\\
William Shakespeare. [\dots] the playwright of this famous tragedy [\dots]
($1564$--$1623$).
&
William Shakespeare. [\dots] a famous English playwright who lived between
approximately $1564$--$1616$ [\dots] including Romeo \& Juliet! \\
\bottomrule
\end{tabular}
\caption{Text continuations (greedy, deterministic; lightly trimmed, [\dots]).
On medical prompts the base stays vague or fails to answer, while the medicine
adapter gives precise clinical content and the register of PubMed (drug
mechanism, ``reducing hepatic gluconeogenesis''; a diagnosis with management). On
a general-knowledge control both answer correctly, i.e.\ the domain adaptation
does not cost general factual recall. Residual chat artifacts (``Hello!\dots'')
appear because the base is partly instruction-tuned.}
\label{tab:med-text}
\end{table}

\begin{figure}[t]
\centering
\setlength{\tabcolsep}{1.5pt}\renewcommand{\arraystretch}{1.0}
\begin{tabular}{@{}cccc@{}}
\rotatebox{90}{\footnotesize\ \ Base} &
\includegraphics[width=0.30\textwidth]{figs/med_comparison/prompt0_base.png} &
\includegraphics[width=0.30\textwidth]{figs/med_comparison/prompt1_base.png} &
\includegraphics[width=0.30\textwidth]{figs/med_comparison/prompt6_base.png} \\
\rotatebox{90}{\footnotesize Med LoRA} &
\includegraphics[width=0.30\textwidth]{figs/med_comparison/prompt0_adapter.png} &
\includegraphics[width=0.30\textwidth]{figs/med_comparison/prompt1_adapter.png} &
\includegraphics[width=0.30\textwidth]{figs/med_comparison/prompt6_adapter.png} \\
\end{tabular}
\caption{Image generation under the text--medicine LoRA: same prompts and seeds,
adapter off (top) vs.\ on (bottom). Generations are essentially unchanged,
confirming the text adapter does not degrade the generation direction (cf.\ the
near-zero Gen-CMMD).}
\label{fig:med-t2i}
\end{figure}


\section{Welsh text-direction qualitative examples}
\label{app:welsh}
Base-vs-adapter continuations for the Welsh text-LoRA ($r{=}32$, $\lambda{=}0.1$);
the base is the identical model with the adapter disabled. Greedy decoding,
repetition penalty $1.3$. Two readings: (i) on \emph{Welsh} seeds the adapter
sustains Welsh while the base breaks into English/assistant text (acquisition);
(ii) on the \emph{English} seed the adapter stays in English (no bleed). As with
poetry, greedy decoding is locally repetitive; sampled decoding ($T{=}0.9$, last
row) is more fluent.

\begin{table}[t]\centering\small
\renewcommand{\arraystretch}{1.3}
\begin{tabular}{@{}p{0.45\linewidth}p{0.45\linewidth}@{}}
\toprule
\multicolumn{2}{@{}l}{\textit{Welsh seed:} ``Mae Cymru yn wlad\,\dots'' (\textit{Wales is a country})}\\
\midrule
\textbf{Base (adapter off)} & \textbf{Welsh LoRA}\\[2pt]
\footnotesize fod â'r eich, a dda i gael \dots\ Hello Mae! I am happy to help you
with your Welsh language question.
&
\footnotesize fwyaf o'r gogledd-ddinasol a ddau 10,578.234 iawn (a chyfansoddwyd
gan yr Aelodau Canolfan Llyfrau Genedlaethol). Mae hynny wedi ei bod arall am\,\dots\\
\midrule
\multicolumn{2}{@{}l}{\textit{English seed:} ``The history of Wales\,\dots''}\\
\midrule
\textbf{Base (adapter off)} & \textbf{Welsh LoRA}\\[2pt]
\footnotesize is a long and complex one, with many different cultures influencing
its development over the centuries. Hello!
&
\footnotesize is a long and complex one, with many different cultures influencing
the country over time. I will be writing about some key events in Welsh History
that have shaped our nation today\,\dots\\
\midrule
\multicolumn{2}{@{}l}{\textit{Welsh seed (sampled, $T{=}0.9$):} ``Yr oedd unwaith\,\dots'' (\textit{There was once})}\\
\midrule
\textbf{Base (adapter off)} & \textbf{Welsh LoRA}\\[2pt]
\footnotesize Hello world \dots\ Wyt yn agor fod dinio? Nochthawd, bye!
&
\footnotesize a byddai'n anhysbys, neu llawn barwas yng ngogleddol rhagolig sydd
wedi bod yn enwu am fod fel cysgu gerfansau pencampell\,\dots\\
\bottomrule
\end{tabular}
\caption{Welsh text-direction qualitative examples ($r{=}32$, $\lambda{=}0.1$). On
Welsh seeds the adapter sustains Welsh where the base slips into English/assistant
text; on the English seed it stays in English (no bleed). Greedy rows are locally
repetitive; the sampled row is more fluent.}
\label{tab:welsh-examples}
\end{table}

\begin{figure}[t]\centering
\setlength{\tabcolsep}{1.5pt}\renewcommand{\arraystretch}{0.8}\small
\begin{tabular}{cc}
\textbf{Base} & \textbf{Welsh LoRA}\\
\includegraphics[width=0.30\linewidth]{figs/t2i_text/welsh/prompt0_base.png} &
\includegraphics[width=0.30\linewidth]{figs/t2i_text/welsh/prompt0_adapter.png}\\
\multicolumn{2}{c}{\scriptsize\textit{a red fox sitting in a snowy forest at dawn}}\\
\addlinespace[1pt]
\includegraphics[width=0.30\linewidth]{figs/t2i_text/welsh/prompt2_base.png} &
\includegraphics[width=0.30\linewidth]{figs/t2i_text/welsh/prompt2_adapter.png}\\
\multicolumn{2}{c}{\scriptsize\textit{a close-up portrait of an elderly fisherman with a wool cap}}\\
\addlinespace[1pt]
\includegraphics[width=0.30\linewidth]{figs/t2i_text/welsh/prompt4_base.png} &
\includegraphics[width=0.30\linewidth]{figs/t2i_text/welsh/prompt4_adapter.png}\\
\multicolumn{2}{c}{\scriptsize\textit{a futuristic city skyline at night with neon lights}}\\
\addlinespace[1pt]
\includegraphics[width=0.30\linewidth]{figs/t2i_text/welsh/prompt6_base.png} &
\includegraphics[width=0.30\linewidth]{figs/t2i_text/welsh/prompt6_adapter.png}\\
\multicolumn{2}{c}{\scriptsize\textit{a plate of spaghetti with tomato sauce and fresh basil}}\\
\addlinespace[1pt]
\includegraphics[width=0.30\linewidth]{figs/t2i_text/welsh/prompt7_base.png} &
\includegraphics[width=0.30\linewidth]{figs/t2i_text/welsh/prompt7_adapter.png}\\
\multicolumn{2}{c}{\scriptsize\textit{an astronaut riding a horse on the surface of Mars}}\\
\end{tabular}
\caption{Image generation under the Welsh text adapter ($r{=}32$, $\lambda{=}0.1$),
base vs adapter on identical prompts and seeds. Generations stay close to the base,
consistent with the near-zero Gen-CMMD ($1.7\times10^{-3}$,
Table~\ref{tab:targets}).}
\label{fig:t2i_welsh}
\end{figure}


\section{Free-form generation samples (MAUVE)}
\label{app:mauve}
The base's low absolute MAUVE-to-WikiText ($0.38$) is not degenerate text: base
continuations are coherent, encyclopedic English (each starts from a fixed
WikiText prefix, then the model continues). MAUVE values are relative to our fixed
protocol ($N{\le}2000$, GPT-2-large features) and are meaningful for the
within-table comparison, not as absolute cross-paper numbers.

\begin{table}[h]\centering\footnotesize
\renewcommand{\arraystretch}{1.2}
\begin{tabular}{@{}p{0.95\linewidth}@{}}
\toprule
\textbf{Base, free-form English continuations (WikiText prefix + generation)}\\
\midrule
\footnotesize Sails persisted on `cruising ironclads' for much longer. During the
1860s, the French navy had produced the Alma and La V\'erendaye, two of the first
ironclads to be built\dots\\
\midrule
\footnotesize Mortier accepted the bait of a rumored Russian retreat. In the early
morning of 11 November he and Gazan departed from D\"urenstein to seize the
abandoned fort\dots\\
\midrule
\footnotesize Following a hiatus of two years, Louis Neimand was hired as rabbi in
1963. Born in New York City in 1919, Neimand grew up in Brooklyn and graduated
from Cornell University in 1941\dots\\
\bottomrule
\end{tabular}
\caption{Base free-form English generations are coherent WikiText-style prose
(cf.\ MAUVE $0.38$ in Table~\ref{tab:mauve}); the low absolute value reflects the
distributional gap of a multimodal assistant model from encyclopedic text, not
degeneracy.}
\label{tab:mauve-samples}
\end{table}


\section{Cross-modal transfer of the acquired language}
\label{app:xmodal}
The probes reported in Table~\ref{tab:xmodal} are constructed as follows. For (A)
we generate the same content from parallel English/Welsh prompts
(Table~\ref{tab:xmodal-welsh-prompts}) and score every image against the same
\emph{English} reference with CLIPScore (a strong English CLIP; the Welsh side is
thus scored on meaning, not on a Welsh-blind text encoder). For (C) the image and
answer format are fixed and only the question language changes; answers are
scored by the Yes/No logit, so output-language drift cannot corrupt them. Because
the base barely knows Welsh, an improvement would be positive cross-modal
transfer of the new language, not a retention effect.

\begin{table}[h]\centering\footnotesize
\setlength{\tabcolsep}{6pt}
\begin{tabular}{@{}p{0.46\linewidth}p{0.46\linewidth}@{}}
\toprule
\textbf{English prompt} & \textbf{Welsh prompt} \\
\midrule
a red fox sitting in a snowy forest & llwynog coch yn eistedd mewn coedwig eira \\
a yellow school bus on a city street & bws ysgol melyn ar stryd y ddinas \\
a slice of chocolate cake on a plate & tafell o gacen siocled ar bl\^at \\
an old stone castle on a hill & castell carreg hen ar fryn \\
a lighthouse on a rocky coast at sunset & goleudy ar arfordir creigiog ar fachlud \\
a plate of fish and chips & pl\^at o bysgod a sglodion \\
\bottomrule
\end{tabular}
\caption{Sample of the parallel English/Welsh prompts for the text$\to$image
probe (Table~\ref{tab:xmodal}, A); the full set has $30$ pairs. The Welsh is
a translation of the English, and both are scored against the same English
reference.}
\label{tab:xmodal-welsh-prompts}
\end{table}

\begin{figure}[h]\centering
\setlength{\tabcolsep}{1.5pt}\renewcommand{\arraystretch}{0.9}\footnotesize
\begin{tabular}{ccc}
\textbf{English $\to$ Base} & \textbf{Welsh $\to$ Base} & \textbf{Welsh $\to$ Welsh LoRA} \\
\includegraphics[width=0.3\linewidth]{figs/xmodal_welsh/p00_s0_base_en.png} &
\includegraphics[width=0.3\linewidth]{figs/xmodal_welsh/p00_s0_base_cy.png} &
\includegraphics[width=0.3\linewidth]{figs/xmodal_welsh/p00_s0_adapter_cy.png} \\
\multicolumn{3}{l}{\scriptsize\textit{a red fox\dots / llwynog coch yn eistedd mewn coedwig eira}} \\
\addlinespace[2pt]
\includegraphics[width=0.3\linewidth]{figs/xmodal_welsh/p02_s0_base_en.png} &
\includegraphics[width=0.3\linewidth]{figs/xmodal_welsh/p02_s0_base_cy.png} &
\includegraphics[width=0.3\linewidth]{figs/xmodal_welsh/p02_s0_adapter_cy.png} \\
\multicolumn{3}{l}{\scriptsize\textit{a yellow school bus\dots / bws ysgol melyn ar stryd y ddinas}} \\
\addlinespace[2pt]
\includegraphics[width=0.3\linewidth]{figs/xmodal_welsh/p02_s0_base_en.png} &
\includegraphics[width=0.3\linewidth]{figs/xmodal_welsh/p02_s0_base_cy.png} &
\includegraphics[width=0.3\linewidth]{figs/xmodal_welsh/p02_s0_adapter_cy.png} \\
\multicolumn{3}{l}{\scriptsize\textit{a yellow school bus\dots / bws ysgol melyn ar stryd y ddinas}} \\
\addlinespace[2pt]
\includegraphics[width=0.3\linewidth]{figs/xmodal_welsh/p04_s0_base_en.png} &
\includegraphics[width=0.3\linewidth]{figs/xmodal_welsh/p04_s0_base_cy.png} &
\includegraphics[width=0.3\linewidth]{figs/xmodal_welsh/p04_s0_adapter_cy.png} \\
\multicolumn{3}{l}{\scriptsize\textit{a slice of chocolate cake\dots / tafell o gacen siocled ar bl\^at}} \\
\end{tabular}
\caption{Text$\to$image under the Welsh adapter. Left: the English prompt renders
correctly (base). Middle and right: the same content requested in Welsh fails for
both base and adapter, i.e.\ the acquired language does not transfer to the
generation pathway (cf.\ Table~\ref{tab:xmodal}).}
\label{fig:xmodal-welsh}
\end{figure}

\end{document}
